
\chapter[定理 7.4.132 中的微分分次李代数同态]
{定理 7.4.132 中的\\\mbox{微分分次李代数同态}}
\label{app:a4}

在定理~\ref{thm:7-4-132} 中，我们定义了同态 $\Phi$。本节证明
$\Phi$ 是微分分次李代数同态。沿用定理~\ref{thm:7-4-132}
证明中的记号
$\widetilde Y,Y(s),Z(s),\mathbf x_i,\mathbf x,x(s),y(s),*_i$ 等。

首先回顾分别作用在下列两个空间上的微分 $d$ 与 $d'$%
\translatornote{源书第二个空间写作
$\operatorname{Der}(\operatorname{Poly}([0,1],BC[1])$，少一个右括号，
且与式~\eqref{eq:7-4-133} 的靶空间写法不一致。译文依该式及下文
以 $\mathbf x_i\in\operatorname{Poly}([0,1],C[1])$ 为输入的型别
采用下式写法；见疑误 A-0724。}：
\[
\begin{aligned}
 d&:\quad
 \operatorname{Poly}([0,1],\operatorname{Der}B(C[1])),\\
 d'&:\quad
 \operatorname{Der}B(\operatorname{Poly}([0,1],C[1])).
\end{aligned}
\]
对
$\widetilde Y\in
\operatorname{Poly}([0,1],\operatorname{Der}B(C[1]))$，有
\[
 d(\widetilde Y)
 =[\overline{\mathfrak m}_1,Y(s)]^{\mathrm{sup}}
 +ds\wedge\left(
   \frac{\partial Y(s)}{\partial s}
   -[\overline{\mathfrak m}_1,Z(s)]^{\mathrm{sup}}
 \right),
\]
其中 $\overline{\mathfrak m}_1$ 是作用在 $C[1]$ 上的
$\mathfrak m_{1,0}$。对
$X\in\operatorname{Der}B(\operatorname{Poly}([0,1],C[1]))$，有
\[
 d'(X)=[\overline{\mathfrak m}_1^{\operatorname{poly}},X]^{\mathrm{sup}},
\]
其中 $\overline{\mathfrak m}_1^{\operatorname{poly}}$ 是作用在
$\operatorname{Poly}([0,1],BC[1])$ 上的 $\mathfrak m_{1,0}$。置
\[
 \overline{\mathfrak m}_1^{\operatorname{poly}}
 =\overline{\mathfrak m}_1^{\operatorname{poly},\circ}
  +ds\wedge\frac{\partial}{\partial s}.
\]
于是
\[
 \overline{\mathfrak m}_1^{\operatorname{poly},\circ}
 \bigl(x(s)+ds\wedge y(s)\bigr)
 =\overline{\mathfrak m}_1(x(s))
  -ds\wedge\overline{\mathfrak m}_1(y(s)).
\]
从而
\[
 d'(X)
 =[\overline{\mathfrak m}_1^{\operatorname{poly},\circ},X]^{\mathrm{sup}}
 +\left[ds\wedge\frac{\partial}{\partial s},X\right]^{\mathrm{sup}}.
\]

另一方面，
\begingroup\small
\begin{equation}
\begin{aligned}
 \Phi(d(\widetilde Y))(\mathbf x_1,\ldots,\mathbf x_k)
 ={}&[\overline{\mathfrak m}_1,Y(s)]^{\mathrm{sup}}
       (x_1(s),\ldots,x_k(s))\\
 &+ds\wedge\sum_i(-1)^{\widetilde{*}_i}
   [\overline{\mathfrak m}_1,Y(s)]^{\mathrm{sup}}
   (x_1(s),\ldots,y_i(s),\ldots,x_k(s))\\
 &+ds\wedge\left\{
   \frac{\partial Y(s)}{\partial s}
   -[\overline{\mathfrak m}_1,Z(s)]^{\mathrm{sup}}
   \right\}(x_1(s),\ldots,x_k(s)).
\end{aligned}
\tag{A4.1}\label{eq:a4-1}
\end{equation}
\endgroup
这里
\[
 \widetilde{*}_i
 =\sum_{j=1}^{i-1}\deg' x_j
  +\deg[\overline{\mathfrak m}_1,Y(s)]^{\mathrm{sup}}
 =\sum_{j=1}^{i-1}\deg' x_j+\deg Y+1
 =*_i+1.
\]

下面证明 $\Phi(d(\widetilde Y))=d'(\Phi(\widetilde Y))$。注意
\[
 \widehat{\overline{\mathfrak m}}_1^{\operatorname{poly},\circ}
 (\mathbf x_1,\ldots,\mathbf x_k)
 =\sum_i(-1)^{a_i}\,
  \mathbf x_1\otimes\cdots\otimes
  \overline{\mathfrak m}_1^{\operatorname{poly},\circ}(\mathbf x_i)
  \otimes\cdots\otimes\mathbf x_k,
\]
其中
\[
 a_i=\sum_{j=1}^{i-1}\deg' x_j.
\]
因此
\[
 \bigl(\Phi(\widetilde Y)\circ
 \widehat{\overline{\mathfrak m}}_1^{\operatorname{poly},\circ}\bigr)
 (\mathbf x_1,\ldots,\mathbf x_k)
 =\bigl(Y(s)\circ\widehat{\overline{\mathfrak m}}_1\bigr)
  (x_1(s),\ldots,x_k(s))+ds\wedge U(s).
\]
这里
\begingroup\footnotesize
\[
\begin{aligned}
 &U(s)-\bigl(Z(s)\circ\widehat{\overline{\mathfrak m}}_1\bigr)
       (x_1(s),\ldots,x_k(s))\\
 &={}
 \sum_i(-1)^{a_i}\Biggl\{
  \sum_{j<i}(-1)^{b_{ij}}
  Y(s)\bigl(x_1(s),\ldots,y_j(s),\ldots,
             \overline{\mathfrak m}_1(x_i(s)),\ldots,x_k(s)\bigr)\\
 &\hspace{35mm}
  +(-1)^{b_{ii}}
  Y(s)\bigl(x_1(s),\ldots,
             \overline{\mathfrak m}_1(y_i(s)),\ldots,x_k(s)\bigr)\\
 &\hspace{35mm}
  +\sum_{j>i}(-1)^{b_{ij}}
  Y(s)\bigl(x_1(s),\ldots,
             \overline{\mathfrak m}_1(x_i(s)),\ldots,
             y_j(s),\ldots,x_k(s)\bigr)
 \Biggr\},
\end{aligned}
\]
\endgroup
其中 $j<i$ 时 $b_{ij}=*_j$，而 $j\geq i$ 时
$b_{ij}=*_j+1$。

另一方面，对每个固定的 $j$，有
\begingroup\small
\[
\begin{aligned}
 &\widehat{\overline{\mathfrak m}}_1
  (x_1(s),\ldots,y_j(s),\ldots,x_k(s))\\
 ={}&\sum_{i<j}(-1)^{c_{ij}}
  \bigl(x_1(s)\otimes\cdots\otimes
        \overline{\mathfrak m}_1(x_i(s))\otimes\cdots\otimes
        y_j(s)\otimes\cdots\otimes x_k(s)\bigr)\\
 &+(-1)^{c_{jj}}
  x_1(s)\otimes\cdots\otimes
  \overline{\mathfrak m}_1(y_j(s))\otimes\cdots\otimes x_k(s)\\
 &+\sum_{j<i}(-1)^{c_{ij}}
  \bigl(x_1(s)\otimes\cdots\otimes y_j(s)\otimes\cdots\otimes
        \overline{\mathfrak m}_1(x_i(s))\otimes\cdots\otimes x_k(s)\bigr),
\end{aligned}
\]
\endgroup
其中 $i\leq j$ 时 $c_{ij}=a_i$，而 $i>j$ 时
$c_{ij}=a_i-1$。\translatornote{源书中间项写作
$(-1)^{c_{ii}}$，但本段固定的是 $j$，该项不在 $i$ 的求和内；
依紧随其后的定义及下式的同余关系应为 $c_{jj}=a_j$。见疑误
A-0725。}
注意
\[
 a_i+b_{ij}+c_{ij}\equiv *_j+1\pmod 2.
\]
所以
\[
\begin{aligned}
 U(s)={}&-\sum_{j=1}^k(-1)^{*_j}
 \bigl(Y(s)\circ\widehat{\overline{\mathfrak m}}_1\bigr)
 (x_1(s),\ldots,y_j(s),\ldots,x_k(s))\\
 &+\bigl(Z(s)\circ\widehat{\overline{\mathfrak m}}_1\bigr)
 (x_1(s),\ldots,x_k(s)).
\end{aligned}
\]

把 $d'(\Phi(\widetilde Y))(\mathbf x_1,\ldots,\mathbf x_k)$ 写成
下列两项之和：
\[
\begin{aligned}
 \mathrm I
 &=[\overline{\mathfrak m}_1^{\operatorname{poly},\circ},
      \Phi(\widetilde Y)]^{\mathrm{sup}}
   (\mathbf x_1,\ldots,\mathbf x_k),\\
 \mathrm {II}
 &=\left[ds\wedge\frac{\partial}{\partial s},
      \Phi(\widetilde Y)\right]^{\mathrm{sup}}
   (\mathbf x_1,\ldots,\mathbf x_k).
\end{aligned}
\]
先计算 $\mathrm I$：
\begingroup\small
\begin{equation}
\begin{aligned}
 \mathrm I
={}&\bigl(\overline{\mathfrak m}_1^{\operatorname{poly},\circ}
       \circ\Phi(\widetilde Y)\bigr)
       (\mathbf x_1,\ldots,\mathbf x_k)\\
 &+(-1)^{\deg Y+1}
   \bigl(\Phi(\widetilde Y)\circ
   \widehat{\overline{\mathfrak m}}_1^{\operatorname{poly},\circ}\bigr)
   (\mathbf x_1,\ldots,\mathbf x_k)\\
={}&(\overline{\mathfrak m}_1\circ Y(s))
    (x_1(s),\ldots,x_k(s))\\
 &+(-1)^{\deg Y+1}
   \bigl(Y(s)\circ\widehat{\overline{\mathfrak m}}_1\bigr)
   (x_1(s),\ldots,x_k(s))\\
 &-\sum_{i=1}^k(-1)^{*_i}ds\wedge
   (\overline{\mathfrak m}_1\circ Y(s))
   (x_1(s),\ldots,y_i(s),\ldots,x_k(s))\\
 &-ds\wedge(\overline{\mathfrak m}_1\circ Z(s))
   (x_1(s),\ldots,x_k(s))\\
 &-(-1)^{\deg Y+1}\sum_{j=1}^k(-1)^{*_j}ds\wedge
   \bigl(Y(s)\circ\widehat{\overline{\mathfrak m}}_1\bigr)
   (x_1(s),\ldots,y_j(s),\ldots,x_k(s))\\
 &+(-1)^{\deg Y+1}ds\wedge
   \bigl(Z(s)\circ\widehat{\overline{\mathfrak m}}_1\bigr)
   (x_1(s),\ldots,x_k(s))\\
={}&[\overline{\mathfrak m}_1,Y(s)]^{\mathrm{sup}}
   (x_1(s),\ldots,x_k(s))\\
 &+ds\wedge\Biggl\{
   \sum_i(-1)^{\widetilde{*}_i}
   [\overline{\mathfrak m}_1,Y(s)]^{\mathrm{sup}}
   (x_1(s),\ldots,y_i(s),\ldots,x_k(s))\\
 &\hspace{37mm}
   -[\overline{\mathfrak m}_1,Z(s)]^{\mathrm{sup}}
   (x_1(s),\ldots,x_k(s))\Biggr\}.
\end{aligned}
\tag{A4.2}\label{eq:a4-2}
\end{equation}
\endgroup

再计算
\[
\begin{aligned}
 \mathrm {II}
={}&\left[ds\wedge\frac{\partial}{\partial s},
      \Phi(\widetilde Y)\right]^{\mathrm{sup}}
      (\mathbf x_1,\ldots,\mathbf x_k)\\
={}&ds\wedge\frac{\partial}{\partial s}
   \bigl(Y(s)(x_1(s),\ldots,x_k(s))\bigr)\\
 &+(-1)^{\deg Y+1}ds\wedge
   \sum_i(-1)^{*_i+d_i}Y(s)
   \left(x_1(s),\ldots,
   \frac{\partial x_i(s)}{\partial s},\ldots,x_k(s)\right),
\end{aligned}
\]
其中
\[
 d_i=\sum_{j=1}^{i-1}\deg' x_j.
\]
注意
$\deg Y+1+*_i+d_i\equiv1\pmod 2$，故
\begin{equation}
 \mathrm {II}
 =ds\wedge\left(\frac{\partial Y(s)}{\partial s}\right)
 (x_1(s),\ldots,x_k(s)).
\tag{A4.3}\label{eq:a4-3}
\end{equation}
比较式~\eqref{eq:a4-1}、\eqref{eq:a4-2} 与~\eqref{eq:a4-3}，得到
\[
 \Phi(d(\widetilde Y))=d'(\Phi(\widetilde Y)).
\]
因此 $\Phi$ 是微分分次模同态。

下面证明 $\Phi$ 是分次李代数同态。设
\[
 \widetilde Y_i=Y_i(s)+ds\wedge Z_i(s),\qquad i=1,2,
\]
是 $\operatorname{Poly}([0,1],\operatorname{Der}B(C[1]))$ 的元素。
于是
\begingroup\small
\[
\begin{aligned}
 [\widetilde Y_1,\widetilde Y_2]^{\mathrm{sup}}
={}&[Y_1(s),Y_2(s)]^{\mathrm{sup}}\\
 &+(-1)^{\deg Y_1}ds\wedge
   \bigl(Y_1(s)\circ\widehat Z_2(s)\bigr)\\
 &+ds\wedge\bigl(Z_1(s)\circ\widehat Y_2(s)\bigr)\\
 &+(-1)^{\deg Y_1\deg Y_2+1}
  \Bigl\{(-1)^{\deg Y_2}ds\wedge
  \bigl(Y_2(s)\circ\widehat Z_1(s)\bigr)\\
 &\hspace{35mm}
  +ds\wedge\bigl(Z_2(s)\circ\widehat Y_1(s)\bigr)\Bigr\}.
\end{aligned}
\]
\endgroup
所以
\begingroup\small
\[
\begin{aligned}
 &\Phi([\widetilde Y_1,\widetilde Y_2]^{\mathrm{sup}})
   (\mathbf x_1,\ldots,\mathbf x_k)\\
={}&[Y_1(s),Y_2(s)]^{\mathrm{sup}}(x_1(s),\ldots,x_k(s))\\
 &+ds\wedge\sum_i(-1)^{\epsilon_i}
   [Y_1(s),Y_2(s)]^{\mathrm{sup}}
   \bigl(x_1(s),\ldots,y_i(s),\\[-2pt]
 &\hspace{54mm}{}\ldots,x_k(s)\bigr)\\
 &+ds\wedge\Bigl\{
   (-1)^{\deg Y_1}\bigl(Y_1(s)\circ\widehat Z_2(s)\bigr)
   +\bigl(Z_1(s)\circ\widehat Y_2(s)\bigr)\\
 &\hspace{29mm}
   +(-1)^{\deg Y_1\deg Y_2+\deg Y_2+1}
    \bigl(Y_2(s)\circ\widehat Z_1(s)\bigr)\\
 &\hspace{29mm}
   +(-1)^{\deg Y_1\deg Y_2+1}
    \bigl(Z_2(s)\circ\widehat Y_1(s)\bigr)
   \Bigr\}(x_1(s),\ldots,x_k(s)),
\end{aligned}
\]
\endgroup
其中
\[
 \epsilon_i=\sum_{k=1}^{i-1}\deg' x_k
 +\deg[Y_1,Y_2]^{\mathrm{sup}}.
\]

另一方面，
\begingroup\small
\[
\begin{aligned}
 &\bigl(\Phi(\widetilde Y_1)\circ
   \widehat{\Phi(\widetilde Y_2)}\bigr)
   (\mathbf x_1,\ldots,\mathbf x_k)\\
={}&\bigl(Y_1(s)\circ\widehat Y_2(s)\bigr)
    (x_1(s),\ldots,x_k(s))\\
 &+\sum_{i=1}^k(-1)^{\xi_i}ds\wedge
   \bigl(Y_1(s)\circ\widehat Y_2(s)\bigr)
   (x_1(s),\ldots,y_i(s),\ldots,x_k(s))\\
 &+(-1)^{\deg Y_1}ds\wedge
   \bigl(Y_1(s)\circ\widehat Z_2(s)\bigr)
   (x_1(s),\ldots,x_k(s))\\
 &+ds\wedge\bigl(Z_1(s)\circ\widehat Y_2(s)\bigr)
   (x_1(s),\ldots,x_k(s)),
\end{aligned}
\]
\endgroup
其中
\[
 \xi_i=\left(\sum_{k=1}^{i-1}\deg' x_k+\deg Y_2\right)+\deg Y_1;
\]
并且
\begingroup\small
\[
\begin{aligned}
 &\bigl(\Phi(\widetilde Y_2)\circ
   \widehat{\Phi(\widetilde Y_1)}\bigr)
   (\mathbf x_1,\ldots,\mathbf x_k)\\
={}&\bigl(Y_2(s)\circ\widehat Y_1(s)\bigr)
    (x_1(s),\ldots,x_k(s))\\
 &+\sum_{i=1}^k(-1)^{\eta_i}ds\wedge
   \bigl(Y_2(s)\circ\widehat Y_1(s)\bigr)
   (x_1(s),\ldots,y_i(s),\ldots,x_k(s))\\
 &+(-1)^{\deg Y_2}ds\wedge
   \bigl(Y_2(s)\circ\widehat Z_1(s)\bigr)
   (x_1(s),\ldots,x_k(s))\\
 &+ds\wedge\bigl(Z_2(s)\circ\widehat Y_1(s)\bigr)
   (x_1(s),\ldots,x_k(s)),
\end{aligned}
\]
\endgroup
其中
\[
 \eta_i=\left(\sum_{k=1}^{i-1}\deg' x_k+\deg Y_1\right)+\deg Y_2.
\]
注意 $\xi_i=\eta_i=\epsilon_i$。由超交换子的定义，
\[
 [\Phi(\widetilde Y_1),\Phi(\widetilde Y_2)]^{\mathrm{sup}}
 =\Phi(\widetilde Y_1)\circ\widehat{\Phi(\widetilde Y_2)}
 +(-1)^{\deg Y_1\deg Y_2+1}
  \Phi(\widetilde Y_2)\circ\widehat{\Phi(\widetilde Y_1)}.
\]
因此
\begingroup\small
\[
\begin{aligned}
 &[\Phi(\widetilde Y_1),\Phi(\widetilde Y_2)]^{\mathrm{sup}}\\
 &\hspace{12mm}{}
   (\mathbf x_1,\ldots,\mathbf x_k)\\
={}&[Y_1(s),Y_2(s)]^{\mathrm{sup}}(x_1(s),\ldots,x_k(s))\\
 &+ds\wedge\sum_i(-1)^{\epsilon_i}
   [Y_1(s),Y_2(s)]^{\mathrm{sup}}
   \bigl(x_1(s),\ldots,y_i(s),\\[-2pt]
 &\hspace{54mm}{}\ldots,x_k(s)\bigr)\\
 &+ds\wedge\Bigl\{
   (-1)^{\deg Y_1}\bigl(Y_1(s)\circ\widehat Z_2(s)\bigr)
   \\[-2pt]
 &\hspace{26mm}{}
   +(-1)^{\deg Y_1\deg Y_2+1+\deg Y_2}
    \bigl(Y_2(s)\circ\widehat Z_1(s)\bigr)
   \Bigr\}\\[-2pt]
 &\hspace{20mm}{}(x_1(s),\ldots,x_k(s))\\
 &+ds\wedge\Bigl\{
   \bigl(Z_1(s)\circ\widehat Y_2(s)\bigr)
   +(-1)^{\deg Y_1\deg Y_2+1}
    \bigl(Z_2(s)\circ\widehat Y_1(s)\bigr)
   \Bigr\}\\[-2pt]
 &\hspace{20mm}{}(x_1(s),\ldots,x_k(s))\\
={}&\Phi([\widetilde Y_1,\widetilde Y_2]^{\mathrm{sup}})\\[-2pt]
 &\hspace{12mm}{}(\mathbf x_1,\ldots,\mathbf x_k).
\end{aligned}
\]
\endgroup
故 $\Phi$ 是分次李代数同态。前面还证明了它是微分分次模同态，
所以 $\Phi$ 是微分分次李代数同态。
